% =====================================================
% Problem Statement
% =====================================================
\section{Problem Statement}
The task of DR grading from fundus photographs is fundamentally an \emph{ordinal} five-class classification problem, where the severity grades follow a natural order ($0 < 1 < 2 < 3 < 4$) and the cost of misclassification grows with the distance between the predicted and true grade. Standard image classification models that treat the grades as independent categorical labels therefore discard valuable structural information and tend to produce clinically implausible errors. The problem is further complicated by a severe class imbalance, since the majority of patients exhibit no or mild disease, leaving the higher-severity grades drastically under-represented in available datasets. The objective of this work is to design a convolutional neural network that respects the ordinal nature of DR severity, mitigates class imbalance, and produces interpretable, clinically aligned predictions evaluated through the Quadratic Weighted Kappa (QWK) metric.

\subsection{Baseline Model}
We approached the problem the way any standard image classification task would be approached: with an ImageNet-pretrained CNN, a linear head re-mapped to the five DR grades, and a vanilla cross-entropy loss. Three backbones of progressively increasing capacity were chosen so that we could observe how depth and architectural sophistication interact with our small ($\sim$3.6k-image) and heavily imbalanced training set. The first model, used as our reference baseline, is \textbf{ResNet-18} ($\sim$11.2M parameters)---a residual network deep enough to learn rich features but light enough to train end-to-end on a single 4\,GB consumer GPU. The second is \textbf{ResNet-50} ($\sim$25M parameters), chosen to test whether simply going deeper improves performance under identical training conditions. The third is \textbf{EfficientNet-B0} ($\sim$5M parameters), selected as a more parameter-efficient, compound-scaled architecture that often dominates the medical imaging literature for fundus tasks. Across all three runs the optimisation recipe was kept deliberately fixed---Adam optimiser at $\text{lr}=10^{-4}$, $224 \times 224$ input resolution, $10$ training epochs, an $80/20$ stratified train--validation split, and unweighted \textbf{categorical cross-entropy} as the loss function---so that any differences in performance can be attributed to the architecture rather than to optimisation tricks.

\subsection{Baseline Results}
The ResNet-18 baseline achieved a best validation Quadratic Weighted Kappa of \textbf{QWK\,$=0.8819$} and accuracy of \textbf{$0.8349$}. The deeper ResNet-50 and EfficientNet-B0 runs produced essentially the same picture (QWK $=0.8678$ and $0.8771$ respectively), confirming that capacity alone is not the bottleneck. However, the training dynamics revealed a much more important problem: \textbf{severe overfitting}. The training loss collapsed from $0.64$ at epoch~1 to $0.022$ by epoch~10, while the validation loss steadily climbed from $0.49$ to $0.73$ over the same window---a textbook divergence between train and validation curves. The per-class report makes the consequence visible: while the dominant Grade~0 class is recognised almost perfectly (F1 $=0.96$), the minority Grade~3 collapses to F1 $=0.49$ and the clinically critical Grade~4 to F1 $=0.60$, with the model effectively memorising the majority class rather than learning lesion-level features. The deeper ResNet-50 fared even worse on Grades~3 and~4 (F1 $=0.37$ and $0.50$), suggesting that adding parameters without addressing data limitations actively hurts the under-represented classes. In short, all three architectures, when trained na\"{i}vely with cross-entropy, produced strong-looking aggregate numbers that masked clinically unacceptable failure on advanced disease.

\subsection{Defining Objectives}
The baseline experiments make clear that \emph{any} typical image classification CNN, when transplanted directly onto APTOS-2019, fails in the same predictable way: it overfits a small, heavily class-imbalanced medical dataset whose visual characteristics (radial fundus geometry, subtle pathology, dataset-specific colour casts and illumination) differ sharply from the natural-image distribution the backbone was pretrained on. The objectives of this project therefore reorient away from chasing yet another architecture and towards directly attacking the root causes:
\begin{enumerate}
    \item \textbf{Reduce overfitting} through medical-aware data augmentation, regularisation (dropout, weight decay), and stronger normalisation of the train/validation curves.
    \item \textbf{Fine-tune ResNet rather than train from scratch}, in particular by freezing the early convolutional stages---which encode low-level edge/texture filters that transfer well from ImageNet---and re-training only the later, task-specific layers and the classification head.
    \item \textbf{Explore alternatives to flat softmax classification} that respect the ordinal structure and class imbalance of DR severity, including weighted cross-entropy, focal loss for hard-example mining, and the CORAL rank-consistent ordinal regression framework.
\end{enumerate}
